% !TEX program = xelatex
\documentclass[12pt,a4paper]{article}

% ── Font & tiếng Việt (XeLaTeX) ──────────────────────────────────────────────
\usepackage{fontspec}
\usepackage{polyglossia}
\setdefaultlanguage{vietnamese}
\setotherlanguage{english}
\setmainfont{Times New Roman}
\setsansfont{Arial}
\setmonofont{Courier New}

% ── Layout & màu sắc ─────────────────────────────────────────────────────────
\usepackage[top=2.5cm, bottom=2.5cm, left=3cm, right=2.5cm]{geometry}
\usepackage{xcolor}
\definecolor{primary}{RGB}{30,80,160}
\definecolor{codebg}{RGB}{248,248,248}

% ── Tiêu đề section ──────────────────────────────────────────────────────────
\usepackage{titlesec}
\titleformat{\section}{\large\bfseries\color{primary}}{\thesection.}{0.5em}{}[\titlerule]
\titleformat{\subsection}{\normalsize\bfseries\color{primary}}{\thesubsection.}{0.5em}{}
\titleformat{\subsubsection}{\normalsize\bfseries}{\thesubsubsection.}{0.5em}{}

% ── Các gói tiện ích ─────────────────────────────────────────────────────────
\usepackage{graphicx}
\usepackage{hyperref}
\usepackage{enumitem}
\usepackage{array}
\usepackage{booktabs}
\usepackage{listings}
\usepackage{fancyhdr}
\usepackage{amsmath}

\hypersetup{colorlinks=true, linkcolor=primary, urlcolor=primary}

% ── Header / footer ──────────────────────────────────────────────────────────
\pagestyle{fancy}
\fancyhf{}
\rfoot{\thepage}
\renewcommand{\headrulewidth}{0.4pt}

% ── Code listing ─────────────────────────────────────────────────────────────
\lstset{
  backgroundcolor=\color{codebg},
  basicstyle=\ttfamily\small,
  breaklines=true,
  frame=single,
  rulecolor=\color{gray!40},
  numbers=left,
  numberstyle=\tiny\color{gray},
  xleftmargin=1em
}

% =============================================================================
\begin{document}
% =============================================================================

\section{Tổng quan Dự án}

HaNoi Go là một nền tảng web phục vụ khách du lịch nước ngoài đến Hà Nội, tập
trung vào bốn vấn đề cốt lõi:

\begin{enumerate}
  \item \textbf{Khám phá địa điểm:} Thiếu nền tảng tổng hợp thông tin địa điểm
        du lịch Hà Nội có bản đồ tích hợp, phù hợp với khách nước ngoài.
  \item \textbf{Lập lịch trình tối ưu:} Khách du lịch thường mất nhiều thời gian
        di chuyển không cần thiết do không biết cách nhóm địa điểm theo khu vực
        và giờ mở cửa.
  \item \textbf{Kết nối hoạt động:} Khách solo khó tìm người cùng tham gia một
        hoạt động cụ thể trong thời gian ngắn tại Hà Nội.
  \item \textbf{Gợi ý lưu trú và ăn uống:} Khách nước ngoài thiếu công cụ gợi ý
        hotel và nhà hàng phù hợp với lịch trình và khu vực đang ở.
        (Module này được triển khai tuỳ theo tiến độ.)
\end{enumerate}

% =============================================================================
\section{Module 0: User \& Security}
% =============================================================================

Module này xử lý toàn bộ vòng đời tài khoản người dùng --- từ đăng ký, đăng
nhập, lấy lại mật khẩu, đến quản trị hệ thống bởi Admin. Đây là nền tảng
bảo mật cho tất cả module còn lại, vì mọi tính năng (lưu lịch trình, tạo
activity, chat nhóm) đều yêu cầu xác thực.

% -------------------------------------------------------------------------
\subsection{Đăng ký tài khoản (Sign Up)}

\subsubsection{Bài toán}

Khách du lịch nước ngoài cần tạo tài khoản nhanh chóng để lưu lịch trình và
tham gia activity. Form đăng ký phải tối giản, rõ ràng và thân thiện với người
dùng quốc tế.

\subsubsection{Luồng hoạt động}

\begin{enumerate}
  \item User truy cập trang \texttt{/signup}, nhập \textbf{email},
        \textbf{username}, và \textbf{mật khẩu}.
  \item Frontend validate inline: email đúng định dạng, username không chứa
        ký tự đặc biệt, mật khẩu tối thiểu 8 ký tự.
  \item Backend kiểm tra email và username chưa tồn tại trong DB
        (\texttt{UNIQUE} constraint).
  \item Mật khẩu được hash bằng \textbf{bcrypt} (cost factor 12) trước khi
        lưu --- không bao giờ lưu plain text.
  \item Hệ thống tạo bản ghi trong bảng \texttt{users} với
        \texttt{role = USER}, \texttt{status = ACTIVE}.
  \item Trả về \textbf{JWT access token} và \textbf{refresh token} ngay sau
        khi đăng ký thành công --- user không cần đăng nhập lại.
  \item Redirect đến trang chủ hoặc trang trước đó (nếu user bị redirect do
        chưa đăng nhập).
\end{enumerate}

\subsubsection{Tính năng UI}

\begin{itemize}
  \item \textbf{Form đăng ký:} Ba trường Email, Username, Password với
        label rõ ràng. Nút ``Sign Up'' vô hiệu hoá khi đang gửi request.
  \item \textbf{Inline validation:} Thông báo lỗi xuất hiện ngay dưới từng
        trường khi blur (rời khỏi ô nhập), không đợi submit.
  \item \textbf{Password strength indicator:} Thanh hiển thị độ mạnh mật
        khẩu (weak / fair / strong) dựa trên độ dài và ký tự đặc biệt.
  \item \textbf{Link chuyển trang:} ``Already have an account? Log in''
        dẫn đến trang Login.
  \item \textbf{Error toasts:} Hiển thị lỗi từ server (email đã tồn tại,
        username đã dùng) dưới dạng toast notification.
\end{itemize}

% -------------------------------------------------------------------------
\subsection{Đăng nhập (Login)}

\subsubsection{Bài toán}

User đã có tài khoản cần đăng nhập để truy cập các tính năng cần xác thực.
Hệ thống dùng \textbf{JWT stateless} --- không lưu session trên server, phù
hợp với kiến trúc NestJS + Next.js tách biệt.

\subsubsection{Luồng hoạt động}

\begin{enumerate}
  \item User nhập \textbf{email} và \textbf{mật khẩu} tại trang
        \texttt{/login}.
  \item Backend xác thực: tìm user theo email, so sánh mật khẩu với hash
        bằng \texttt{bcrypt.compare()}.
  \item Kiểm tra \texttt{status}: nếu \texttt{BANNED}, trả về lỗi
        \texttt{403} kèm thông báo tài khoản bị khoá.
  \item Nếu hợp lệ, sinh \textbf{Access Token} (JWT, TTL 15 phút) và
        \textbf{Refresh Token} (JWT, TTL 7 ngày, lưu trong
        \texttt{httpOnly cookie}).
  \item Frontend lưu Access Token trong memory (không lưu localStorage để
        tránh XSS). Refresh Token được gửi tự động qua cookie.
  \item Khi Access Token hết hạn, frontend tự động gọi endpoint
        \texttt{POST /auth/refresh} để lấy token mới mà không cần user
        đăng nhập lại.
\end{enumerate}

\subsubsection{Tính năng UI}

\begin{itemize}
  \item \textbf{Form đăng nhập:} Hai trường Email và Password.
        Nút ``Log In'' có loading spinner khi đang xử lý.
  \item \textbf{Show / hide password:} Icon con mắt để toggle hiển thị
        mật khẩu.
  \item \textbf{``Forgot password?'':} Link dẫn đến trang Reset Password.
  \item \textbf{Redirect sau login:} Nếu user bị chặn do chưa đăng nhập
        (ví dụ cố vào \texttt{/trips}), sau khi login thành công sẽ
        redirect về trang đó.
  \item \textbf{Error state:} Thông báo ``Invalid email or password'' khi
        sai thông tin, không tiết lộ email có tồn tại hay không (tránh
        user enumeration).
\end{itemize}

% -------------------------------------------------------------------------
\subsection{Đặt lại mật khẩu (Reset Password)}

\subsubsection{Bài toán}

User quên mật khẩu cần có cơ chế lấy lại tài khoản an toàn qua email.
Token reset có thời hạn ngắn và chỉ dùng được một lần.

\subsubsection{Luồng hoạt động}

\begin{enumerate}
  \item \textbf{Bước 1 --- Yêu cầu reset:} User vào trang
        \texttt{/forgot-password}, nhập địa chỉ email.
  \item Backend kiểm tra email có tồn tại trong DB. Nếu không tồn tại,
        vẫn trả về thông báo thành công (tránh lộ thông tin tài khoản).
  \item Hệ thống sinh \textbf{reset token} (random 32 bytes, encode
        base64url), lưu hash của token vào bảng \texttt{password\_resets}
        kèm \texttt{expires\_at = now() + 1 giờ}.
  \item Gửi email chứa link dạng
        \texttt{/reset-password?token=<token>} qua \textbf{Nodemailer}
        (SMTP hoặc SendGrid).
  \item \textbf{Bước 2 --- Đặt mật khẩu mới:} User click link trong
        email, được dẫn đến trang \texttt{/reset-password}.
  \item Backend validate token: kiểm tra hash trùng khớp và chưa hết hạn.
        Nếu hợp lệ, cho phép user nhập mật khẩu mới.
  \item Mật khẩu mới được hash bcrypt và cập nhật vào bảng \texttt{users}.
        Token bị xoá khỏi \texttt{password\_resets} ngay sau khi dùng
        (one-time use).
\end{enumerate}

\subsubsection{Schema bổ sung}

\begin{lstlisting}[language=SQL]
CREATE TABLE password_resets (
  id         uuid PRIMARY KEY DEFAULT gen_random_uuid(),
  user_id    uuid REFERENCES users(id) ON DELETE CASCADE,
  token_hash varchar NOT NULL,   -- bcrypt hash cua reset token
  expires_at timestamptz NOT NULL,
  created_at timestamptz DEFAULT now()
);
\end{lstlisting}

\subsubsection{Tính năng UI}

\begin{itemize}
  \item \textbf{Trang Forgot Password (\texttt{/forgot-password}):}
        Ô nhập email, nút ``Send Reset Link''. Sau khi gửi, hiển thị
        thông báo ``Check your inbox'' dù email có tồn tại hay không.
  \item \textbf{Trang Reset Password (\texttt{/reset-password}):}
        Hai ô nhập ``New Password'' và ``Confirm Password'' với kiểm tra
        khớp nhau. Nếu token hết hạn hoặc không hợp lệ, hiển thị thông
        báo lỗi và nút ``Request a new link''.
  \item \textbf{Thành công:} Sau khi đặt lại mật khẩu, redirect về trang
        Login với toast ``Password updated. Please log in.''.
\end{itemize}

% -------------------------------------------------------------------------
\subsection{Admin Site --- Quản lý người dùng}

\subsubsection{Bài toán}

Hệ thống cần giao diện quản trị để Admin có thể xem toàn bộ danh sách
tài khoản, tìm kiếm, và khoá/mở khoá tài khoản vi phạm. Admin Site chỉ
truy cập được bởi tài khoản có \texttt{role = ADMIN} --- mọi endpoint đều
được bảo vệ bởi \textbf{RolesGuard}.

\subsubsection{Phân quyền (RBAC)}

Hệ thống định nghĩa hai role:

\begin{center}
\begin{tabular}{lll}
  \toprule
  \textbf{Role} & \textbf{Quyền} & \textbf{Ghi chú} \\
  \midrule
  \texttt{USER}  & Dùng toàn bộ tính năng app & Mặc định khi đăng ký \\
  \texttt{ADMIN} & USER + quản lý tài khoản   & Gán thủ công trong DB \\
  \bottomrule
\end{tabular}
\end{center}

Trường \texttt{status} kiểm soát khả năng đăng nhập:

\begin{center}
\begin{tabular}{ll}
  \toprule
  \textbf{Status} & \textbf{Ý nghĩa} \\
  \midrule
  \texttt{ACTIVE} & Tài khoản hoạt động bình thường \\
  \texttt{BANNED} & Bị khoá, trả về \texttt{403} khi cố đăng nhập \\
  \bottomrule
\end{tabular}
\end{center}

\subsubsection{Tính năng Admin}

\begin{itemize}
  \item \textbf{Danh sách user (\texttt{GET /admin/users}):} Bảng hiển
        thị toàn bộ tài khoản với các cột: Avatar, Username, Email, Role,
        Status, Ngày tạo. Hỗ trợ \textbf{phân trang} (20 user/trang) và
        \textbf{sắp xếp} theo ngày tạo hoặc username.
  \item \textbf{Tìm kiếm:} Thanh search lọc theo email hoặc username
        (tìm kiếm không phân biệt hoa thường, \texttt{ILIKE}).
  \item \textbf{Khoá tài khoản (\texttt{PATCH /admin/users/:id/ban}):}
        Chuyển \texttt{status} từ \texttt{ACTIVE} sang \texttt{BANNED}.
        User bị đăng xuất ngay lập tức (invalidate token bằng cách tăng
        \texttt{token\_version} trong DB).
  \item \textbf{Mở khoá tài khoản (\texttt{PATCH /admin/users/:id/unban}):}
        Chuyển \texttt{status} về \texttt{ACTIVE}.
  \item \textbf{Xem chi tiết user:} Trang profile đọc-only hiển thị
        thông tin tài khoản và lịch sử hoạt động (số trip đã tạo, số
        activity đã tham gia).
\end{itemize}

\subsubsection{Bảo vệ Route}

\begin{lstlisting}[language={}]
-- Backend: RolesGuard (NestJS)
@Roles('ADMIN')
@UseGuards(JwtAuthGuard, RolesGuard)
@Get('/admin/users')
listUsers() { ... }

-- Frontend: Middleware Next.js
-- Redirect ve /404 neu user.role !== 'ADMIN'
export function middleware(request: NextRequest) {
  const token = request.cookies.get('access_token');
  const payload = decodeJwt(token);
  if (payload.role !== 'ADMIN') {
    return NextResponse.redirect('/404');
  }
}
\end{lstlisting}

\subsubsection{Tính năng UI (Admin)}

\begin{itemize}
  \item \textbf{Layout riêng:} Admin Site có sidebar navigation riêng
        tách biệt với giao diện người dùng thông thường. Truy cập qua
        \texttt{/admin}.
  \item \textbf{User table:} Bảng dữ liệu với cột Status hiển thị badge
        màu (xanh: Active, đỏ: Banned). Mỗi row có dropdown action ``Ban''
        / ``Unban'' với confirm dialog trước khi thực thi.
  \item \textbf{Search bar:} Lọc theo email hoặc username realtime với
        debounce 300ms.
  \item \textbf{Phân trang:} Pagination component ở cuối bảng, hiển thị
        tổng số user và trang hiện tại.
  \item \textbf{Toast notification:} Xác nhận hành động thành công hoặc
        thất bại sau mỗi thao tác Ban/Unban.
\end{itemize}

% =============================================================================
\section{Các Module Chính}
% =============================================================================

\subsection{Khám phá Địa điểm}

\subsubsection{Bài toán}

Khách du lịch nước ngoài cần một nơi tra cứu thông tin địa điểm du lịch tại
Hà Nội một cách trực quan, có bản đồ tích hợp và khả năng lọc theo khu vực
hoặc loại hình. Module này đóng vai trò \textbf{foundation data} cho toàn bộ
hệ thống --- bộ dữ liệu 65 địa điểm du lịch đã được xác thực, sẵn sàng hỗ trợ
tìm đường thời gian thực.

\subsubsection{Tính năng}

\begin{itemize}
  \item \textbf{Danh sách địa điểm:} Hiển thị tên, ảnh, loại hình, khu vực,
        giờ mở cửa, thời gian tham quan ước tính.
  \item \textbf{Filter:} Lọc theo khu vực (Hoàn Kiếm, Ba Đình, Tây Hồ, Đống Đa,
        Hai Bà Trưng, Cầu Giấy) và loại hình (di tích, bảo tàng, hồ, chùa,
        công viên).
  \item \textbf{Map view:} Hiển thị toàn bộ địa điểm dưới dạng pin trên bản đồ
        Leaflet. Click vào pin để xem thông tin chi tiết.
  \item \textbf{Trang chi tiết:} Thông tin đầy đủ kèm vị trí trên bản đồ và
        địa điểm lân cận trong bán kính 1 km (tái dụng PostGIS).
\end{itemize}

\subsubsection{Triển khai}

Dữ liệu địa điểm được lưu trong PostgreSQL với PostGIS. Thay vì lưu giờ mở cửa
dưới dạng JSONB thuần (khó truy vấn và không thể lọc trực tiếp bằng SQL), hệ
thống tách thành các cột riêng biệt giúp thuật toán đọc nhanh và truy vấn hiệu
quả:

\begin{lstlisting}[language=SQL]
CREATE TABLE places (
  id                uuid PRIMARY KEY DEFAULT gen_random_uuid(),
  name              text NOT NULL,
  category          text NOT NULL,
  district          text NOT NULL,
  location          geometry(Point, 4326),
  visit_duration_min int DEFAULT 60,
  -- Lich hoat dong (parsed tu opening hours)
  always_open       boolean DEFAULT false,
  open_days         integer[],  -- [0..6] = [CN,T2..T7]
  open_time_start   time,       -- gio mo cua (VD: 08:00)
  open_time_end     time,       -- gio dong cua (VD: 17:00)
  has_break         boolean DEFAULT false,
  break_start       time,       -- nghi trua bat dau
  break_end         time        -- nghi trua ket thuc
);
\end{lstlisting}

Ví dụ dữ liệu thực tế:

\begin{lstlisting}[language=SQL]
-- Ho Hoan Kiem (mo 24/7)
INSERT INTO places (...) VALUES (
  ..., true, ARRAY[0,1,2,3,4,5,6],
  '00:00', '23:59', false, null, null
);
-- Van Mieu (dong T2, nghi trua)
INSERT INTO places (...) VALUES (
  ..., false, ARRAY[0,2,3,4,5,6],
  '08:00', '17:00', true, '12:00', '13:30'
);
\end{lstlisting}

Truy vấn địa điểm lân cận tái dụng \texttt{ST\_DWithin} --- cùng cơ chế với
Geospatial Search nhưng phục vụ trang chi tiết thay vì tìm kiếm theo GPS
của user.

\begin{lstlisting}
SELECT id, name, category, district,
       always_open, open_days, open_time_start, open_time_end,
       visit_duration_min,
       ST_AsGeoJSON(location) AS geojson
FROM   places
WHERE  district = :district        -- filter khu vuc
  AND  category = ANY(:categories) -- filter loai hinh
ORDER  BY name;
\end{lstlisting}

% -------------------------------------------------------------------------
\subsection{Trip Planner}

\subsubsection{Bài toán}

Người dùng chọn danh sách địa điểm du lịch muốn đến tại Hà Nội cùng số ngày
và điểm xuất phát. Hệ thống tự động sắp xếp lịch trình theo ngày sao cho tổng
quãng đường di chuyển là nhỏ nhất, đồng thời thoả mãn ràng buộc về giờ mở cửa
và thời gian tham quan. Lịch trình được hiển thị trực tiếp trên bản đồ Leaflet.

\subsubsection{Thuật toán}

Bài toán được mô hình hoá dưới dạng một biến thể của \textbf{Travelling Salesman Problem with Time Windows (TSPTW)} có bổ sung ràng buộc đa ngày và khung giờ nghỉ. Hệ thống sử dụng phương pháp \textbf{Cluster-first, Route-second} kết hợp cùng \textbf{Greedy Nearest Neighbor (GNN)}.

\paragraph{Mô hình hoá dữ liệu (Problem Formulation):}
Đầu vào của thuật toán bao gồm tập hợp các địa điểm $P = \{p_1, p_2, \dots, p_n\}$, số ngày tham quan $D$, quỹ thời gian mỗi ngày $[T_{start}, T_{end}]$, và khung giờ nghỉ trưa $[L_{start}, L_{end}]$. Mỗi địa điểm $p_i$ được định nghĩa bởi tọa độ $(lat_i, lng_i)$, thời lượng tham quan $v_i$, và khung giờ mở cửa thực tế trong ngày $[O_{start}(i), O_{end}(i)]$.

\paragraph{Giai đoạn 1: Tiền xử lý và Phân cụm (Clustering)}
Hệ thống loại bỏ các địa điểm đóng cửa hoàn toàn dựa trên mảng \texttt{open\_days}. Tiếp theo, thuật toán K-Means chia tập $P$ thành $D$ cụm $C_1, C_2, \dots, C_D$, mỗi cụm đại diện cho 1 ngày.
\begin{itemize}
    \item \textbf{Khoảng cách (Distance Metric):} Sử dụng công thức Haversine $\text{Dist}(p_i, p_j)$ để gom cụm địa lý.
    \item \textbf{Tính xác định (Determinism):} Việc chọn các tâm cụm (centroids) ban đầu bằng hàm lấy số ngẫu nhiên mặc định (\texttt{Math.random}) gây ra sự thiếu ổn định. Hệ thống đã thay thế bằng một bộ sinh số giả ngẫu nhiên có hạt giống \textbf{(Seeded PRNG)}, trong đó hạt giống $S$ được tính toán dựa trên toạ độ $(lat, lng)$ của các điểm. Điều này đảm bảo thuật toán là một \textit{hàm xác định (deterministic function)} --- cùng một input luôn cho ra một kết quả gom cụm đồng nhất.
    \item \textbf{Cân bằng ngày mở cửa (\texttt{postClusterOpenDaySwap}):} Nếu một điểm $p_i$ bị xếp vào cụm $C_d$ mà ngày $d$ điểm này đóng cửa, thuật toán tự động duyệt tìm cụm $C_{k}$ khả dĩ nhất và hoán đổi để tối ưu khả năng tham quan.
\end{itemize}

\paragraph{Giai đoạn 2: Định tuyến có ràng buộc thời gian (Time-Window GNN)}
Trong mỗi cụm $C_d$, bài toán định tuyến được giải quyết theo hướng tiếp cận tham lam (Greedy). Tại mỗi bước chuyển từ điểm hiện tại (current) sang điểm dự kiến $j$, thời gian thực được lấy qua \textbf{Goong Distance Matrix API}, ký hiệu thời gian di chuyển là $\Delta T(current, j)$.

Các công thức tính toán thời gian cho một ứng viên $j$ được diễn giải chi tiết như sau:
\begin{enumerate}
    \item \textbf{Tính thời gian đến ($T_{arrival}$):}
    \begin{equation}
        T_{arrival}(j) = T_{depart}(current) + \Delta T(current, j)
    \end{equation}

    \item \textbf{Tính thời gian chờ mở cửa ($wait\_open$):} Nếu đến sớm hơn giờ mở cửa.
    \begin{equation}
        wait\_open = \max(0, \; O_{start}(j) - T_{arrival}(j))
    \end{equation}

    \item \textbf{Tính thời gian chờ do nghỉ trưa ($wait\_lunch$):} Gọi thời điểm dự kiến bắt đầu tham quan là $T'_{visit} = T_{arrival}(j) + wait\_open$. Nếu khoảng thời gian tham quan vắt ngang qua giờ nghỉ trưa ($T'_{visit} < L_{end}$ và $T'_{visit} + v_j > L_{start}$), thời gian bắt đầu sẽ bị đẩy xuống cuối giờ nghỉ:
    \begin{equation}
        wait\_lunch = L_{end} - T'_{visit}
    \end{equation}
    Tổng thời gian chờ: $T_{wait}(j) = wait\_open + wait\_lunch$.
    
    \item \textbf{Kiểm tra tính hợp lệ (Feasibility Validation):} Điểm $j$ được coi là hợp lệ để xếp vào ngày nếu thoả mãn:
    \begin{equation}
        T_{arrival}(j) \le O_{end}(j) - v_j \quad \text{(Không đến quá muộn trước giờ đóng cửa)}
    \end{equation}
    \begin{equation}
        T_{arrival}(j) + T_{wait}(j) + v_j \le T_{end} \quad \text{(Kịp hoàn thành trước giờ kết thúc ngày)}
    \end{equation}
    Các điểm không hợp lệ bị đẩy vào danh sách chờ xử lý sau ($P_{dropped}$).

    \item \textbf{Hàm chi phí (Cost Function) và xử lý Tie-Breaking:} 
    Trong số các ứng viên $j$ hợp lệ, thuật toán chọn điểm có $Cost(j)$ nhỏ nhất. Nếu chỉ sử dụng tổng thời gian, hiện tượng "tie-breaking" sẽ xảy ra (ví dụ: các điểm đều phải đợi đến 13:00 mới hết nghỉ trưa sẽ có chi phí bằng nhau, làm mất đi tính liên kết địa lý). Hệ thống thiết kế hàm phạt trọng số:
    \begin{equation}
        Cost(j) = \Big(\Delta T(current, j) \times 2\Big) + \Big(T_{wait}(j) \times 60\Big)
    \end{equation}
    Việc nhân hệ số 2 cho thời gian di chuyển ($\Delta T$ tính bằng giây, $T_{wait}$ tính bằng phút) buộc thuật toán phải ưu tiên tuyệt đối sự liền kề về mặt địa lý, loại bỏ tình trạng lộ trình zic-zắc.
\end{enumerate}

\paragraph{Giai đoạn 3: Giải quyết xung đột (Conflict Resolution - Gap Insertion)}
Với các địa điểm bị loại ($p \in P_{dropped}$), hệ thống tiến hành thuật toán chèn khuyết (Gap Insertion).
Để chèn điểm $p$ vào vị trí thứ $k$ (giữa điểm $k-1$ và điểm $k$ của một ngày $d$), hệ thống tính toán lại thời gian đến và đi:
\begin{align}
    T_{arrival}(p) &= T_{depart}(k-1) + \Delta T(k-1, p) \\
    T_{depart}(p) &= T_{arrival}(p) + T_{wait}(p) + v_p
\end{align}
Hệ thống thực hiện kiểm tra dây chuyền (Cascade Validation) với 3 điều kiện:
\begin{enumerate}
    \item $T_{depart}(p) \le O_{end}(p)$ (Kịp hoàn thành trước khi $p$ đóng cửa).
    \item $T_{depart}(p) \le T_{end}$ (Không vi phạm thời lượng của toàn bộ ngày).
    \item $T_{depart}(p) + \Delta T(p, k) \le T_{arrival}(k)$ (Không làm trễ lịch của điểm đến $k$ tiếp theo).
\end{enumerate}
Nếu thỏa mãn toàn bộ, điểm $p$ được nội suy vào chuỗi. Nếu thử mọi vị trí ở tất cả các ngày đều thất bại, $p$ bị loại hoàn toàn, được đưa vào danh sách \texttt{unscheduled} và trả về cho Frontend hiển thị thông báo "Không thể sắp xếp (Quá giờ/Quá xa)".

\subsubsection{Tính toán khoảng cách}

\textbf{Bước 1 --- Haversine Formula} (K-Means clustering):

\[
  d = 2r \arcsin\!\left(\sqrt{
    \sin^2\!\left(\frac{\Delta\phi}{2}\right)
    + \cos\phi_1\cos\phi_2\sin^2\!\left(\frac{\Delta\lambda}{2}\right)
  }\right)
\]

trong đó $\phi$ là vĩ độ, $\lambda$ là kinh độ, $r = 6{,}371\text{ km}$.

\medskip
\textbf{Bước 2 --- Goong Distance Matrix API} (Greedy Nearest Neighbor routing):

Hệ thống sử dụng \textbf{Goong Distance Matrix API} để định tuyến chính xác thay vì khoảng cách đường chim bay. Đặc biệt, API này được cấu hình với tham số \texttt{vehicle=bike} để phản ánh đúng thời gian di chuyển bằng xe máy tại Hà Nội --- phương tiện di chuyển phổ biến nhất của du khách và người dân địa phương.

Một request duy nhất trả về toàn bộ ma trận thời gian thực tế giữa $N$ địa điểm:

\begin{lstlisting}[language={}]
GET https://rsapi.goong.io/DistanceMatrix?
  origins=21.0285,105.8412|21.0277,105.8412
  &destinations=21.0285,105.8412|21.0277,105.8412
  &vehicle=bike

Response:
  rows: [
    { elements: [{ duration: { value: 0 } }, { duration: { value: 180 } }] },
    { elements: [{ duration: { value: 180 } }, { duration: { value: 0 } }] }
  ] -- giay
\end{lstlisting}

\textbf{Dự phòng lỗi (Fallback \& Rate Limiting):}

Trong quá trình gọi API, nếu gặp lỗi giới hạn tốc độ (\texttt{OVER\_RATE\_LIMIT}), hệ thống sẽ tự động tạm dừng 3 giây và thử lại. Nếu API key gặp sự cố hoặc cạn kiệt request, hệ thống chuyển đổi mượt mà sang công thức Haversine (tốc độ giả định $15\text{ km/h}$) để đảm bảo trải nghiệm người dùng không bị gián đoạn.

\subsubsection{Ví dụ đầu vào / đầu ra}

\begin{lstlisting}[language={}]
Input:
  Dia diem: [Van Mieu, Ho Tay, Hoan Kiem, Bao tang HN, Lang Chu Tich]
  So ngay: 2
  Diem xuat phat: Hoan Kiem

Output:
  Ngay 1: Hoan Kiem (8:00) -> Bao tang HN (10:00) -> Van Mieu (14:00)
  Ngay 2: Lang Chu Tich (8:00) -> Ho Tay (10:30)
\end{lstlisting}

% -------------------------------------------------------------------------
\subsection{Activity \& Group Chat}

\subsubsection{Bài toán}

Khách du lịch solo muốn tìm người cùng tham gia một hoạt động cụ thể tại Hà
Nội trong thời gian ngắn --- nhưng các nền tảng hiện tại (Facebook, Meetup)
không gắn hoạt động trực tiếp lên bản đồ và không hỗ trợ group chat tức thì
sau khi ghép nhóm.

\subsubsection{Luồng hoạt động}

\begin{enumerate}
  \item \textbf{Tạo activity:} User điền tên hoạt động, thời gian, mô tả, số
        người tối đa. Địa điểm có thể chọn từ dataset hoặc nhập tự do (kèm
        toạ độ GPS).
  \item \textbf{Hiển thị trên map:} Activity xuất hiện dưới dạng pin trên bản
        đồ Leaflet. Các user khác click vào pin để xem chi tiết và gửi request
        tham gia.
  \item \textbf{Duyệt thành viên:} Chủ activity nhận thông báo, xem profile
        người xin tham gia và duyệt hoặc từ chối.
  \item \textbf{Group chat:} Sau khi được duyệt, thành viên tự động vào group
        chat của activity. Chat đồng bộ realtime qua WebSocket và Redis Pub/Sub.
\end{enumerate}

\subsubsection{Triển khai}

\begin{lstlisting}[language=SQL]
CREATE TABLE activities (
  id           uuid PRIMARY KEY DEFAULT gen_random_uuid(),
  host_id      uuid REFERENCES users(id),
  place_id     uuid REFERENCES places(id),  -- null neu tu nhap
  title        text NOT NULL,
  description  text,
  location     geometry(Point, 4326),       -- toa do hien tren map
  scheduled_at timestamptz NOT NULL,
  max_members  int DEFAULT 10,
  created_at   timestamptz DEFAULT now()
);

CREATE TABLE activity_members (
  activity_id  uuid REFERENCES activities(id),
  user_id      uuid REFERENCES users(id),
  status       text CHECK (status IN ('pending','approved','rejected')),
  PRIMARY KEY (activity_id, user_id)
);
\end{lstlisting}

\textbf{Realtime:} Tin nhắn group chat được phân phối qua \textbf{WebSocket
Gateway} (NestJS) và \textbf{Redis Pub/Sub}. Mỗi activity có một Redis channel
riêng theo \texttt{activity:\{id\}}. Khi thành viên được duyệt, server tự động
subscribe user vào channel tương ứng.

\begin{lstlisting}[language={}]
-- Phan phoi tin nhan group chat
Channel: activity:{activity_id}
Payload: { userId, username, message, timestamp }

-- Server subscribe user sau khi duyet
redis.subscribe(`activity:${activityId}`)
\end{lstlisting}

% -------------------------------------------------------------------------
\subsection{AI Mentor}

\subsubsection{Bài toán}

Khách du lịch cần một trợ lý ảo am hiểu địa phương để gợi ý địa điểm, giải đáp
thắc mắc về văn hóa và hỗ trợ tối ưu lịch trình dựa trên vị trí hiện tại.

\subsubsection{Triển khai}

Module sử dụng mô hình \textbf{Gemini 2.5 Flash} nhờ khả năng xử lý context lớn
và tốc độ phản hồi nhanh. Hệ thống cung cấp context cho AI bao gồm danh sách
địa điểm lân cận truy vấn từ PostgreSQL/PostGIS và vị trí GPS của người dùng.

\begin{lstlisting}[language={}]
System:
  You are "HanoiGO AI Assistant". 
  Respond in a professional, friendly style with Emojis.
  Prioritize landmarks within 2km if user location is provided.
  Use bold for place names. Respond in Vietnamese.

Prompt:
  Context: [List of 50 nearest landmarks with coordinates]
  User Location: (curLat, curLng)
  Question: "Recommend some points of interest near me."
\end{lstlisting}

% =============================================================================
\section{Kiến trúc Hệ thống}
% =============================================================================

\subsection{Tổng quan}

\begin{center}
\begin{tabular}{ll}
  \toprule
  \textbf{Tầng} & \textbf{Công nghệ} \\
  \midrule
  Frontend         & Next.js 14, Tailwind CSS, Leaflet.js \\
  Backend API      & NestJS (Node.js), WebSocket Gateway  \\
  AI Integration   & Gemini 2.5 Flash (Google AI SDK)      \\
  Routing Engine   & Goong API / OSRM (self-hosted)       \\
  Database Chính   & PostgreSQL + PostGIS                 \\
  Cache / Realtime & Redis Pub/Sub                        \\
  Maps             & CartoDB Voyager + Leaflet            \\
  \bottomrule
\end{tabular}
\end{center}

\subsection{Luồng dữ liệu chính}

\begin{enumerate}
  \item \textbf{Khám phá địa điểm:} Frontend gửi filter → NestJS truy vấn
        PostgreSQL → trả về danh sách + GeoJSON để render Leaflet pins.

  \item \textbf{Trip Planner:} Frontend gửi danh sách địa điểm và số ngày →
        NestJS tiền lọc địa điểm theo \texttt{open\_days} (Bước~0) → K-Means
        clustering (Deterministic với Seeded PRNG) → gọi Goong Distance Matrix API lấy ma trận khoảng cách
        thực tế → GNN Routing có tích hợp hàm chi phí ưu tiên di chuyển và xử lý giờ nghỉ trưa
        ($Cost = TravelTime \times 2 + T_{wait} \times 60$) → chạy Conflict Resolution để điền vào các khe trống (Gap Insertion) → trả về lịch trình JSON + GeoJSON
        kèm danh sách \texttt{infeasible} và \texttt{unscheduled} để cảnh báo người dùng.

  \item \textbf{Activity \& Group Chat:} User tạo activity → lưu DB, render
        pin trên map → user khác request tham gia → chủ duyệt → NestJS
        subscribe user vào Redis channel → WebSocket đồng bộ group chat
        realtime.

  \item \textbf{AI Mentor:} Frontend gửi context lịch trình → NestJS gọi
        Claude API → Claude web search nếu cần → trả về gợi ý dạng chat.
\end{enumerate}

% =============================================================================
\section{Dữ liệu \& Nội dung}
% =============================================================================

Dự án sử dụng \textbf{một bộ dữ liệu duy nhất} (65 địa điểm du lịch trọng điểm)
đã được xác thực tọa độ và thông tin, phục vụ cho cả ba module chính, tránh phân
tán công sức thu thập và bảo trì. Dữ liệu được tổng hợp từ:

\begin{itemize}
  \item \textbf{OpenStreetMap (Overpass API)} --- toạ độ, tên, phân loại
        (giấy phép ODbL, miễn phí). Dữ liệu OSM Việt Nam cũng dùng để khởi
        tạo OSRM routing engine.
  \item \textbf{Bổ sung thủ công} --- giờ mở cửa, thời gian tham quan ước
        tính, ảnh đại diện, mô tả ngắn bằng tiếng Anh.
\end{itemize}

\textbf{Phạm vi dữ liệu:} Chỉ bao gồm địa điểm du lịch (di tích, bảo tàng,
hồ, công viên, chùa,...). Hotel và nhà hàng không nằm trong dataset tĩnh ---
được tra cứu realtime bởi AI Mentor khi cần.

\medskip
Phân loại theo \textbf{khu vực} phục vụ clustering và filter:
Hoàn Kiếm, Ba Đình, Tây Hồ, Đống Đa, Hai Bà Trưng, Cầu Giấy.

\end{document}